\documentclass[12pt]{article}
\usepackage[utf8]{inputenc}
\usepackage[T1]{fontenc}
\usepackage{hyperref}
\title{Pholo Language Resources for K-12 Curriculum Development}
\author{Compiled by Claude Sonnet 5.5 (claude-sonnet-5-5) for Andr\'as Kornai}
\date{}
\begin{document}
\maketitle

\section{Introduction}
Pholo (ISO 639-3 \texttt{yip}, Glottocode phol1235) is a Loloish (Yi) language of Yunnan, China, in the Sino-Tibetan family. Sources give speaker counts of about 30,000 (Wikipedia, 2011) and 35,000 (Joshua Project). Mandarin is increasingly used, especially by younger people. Pholo is not to be confused with Phola (Phu La), a separate language. Web searches found no school materials in Pholo.

\section{Primary Educational Resources}
No primers, textbooks or school programmes for Pholo turned up in the searches. Nothing indicates the language is taught in schools.

\section{Linguistic and Reference Materials}
\begin{itemize}
\item Pholo language, Wikipedia overview: \url{https://en.wikipedia.org/wiki/Pholo_language}
\item Phola language, Wikipedia (a different language, listed here to avoid confusion): \url{https://en.wikipedia.org/wiki/Phola_language}
\end{itemize}
No grammar or dictionary was found in the search results.

\section{Religious and Community Resources}
\begin{itemize}
\item Joshua Project, Pholo in China (demographic and community profile): \url{https://joshuaproject.net/people_groups/21077}
\end{itemize}
No scripture translation was found.

\section{Digital and Audio Resources}
No audio, video or corpus resources for Pholo were found.

\section{Curriculum Development Context}
Very little is documented about Pholo. Any curriculum would need fresh documentation work with the community, starting with an orthography, since none was found. Mandarin is the likely language of instruction at present. Subgrouping within Loloish is unresolved.
\end{document}